\scanpage{23}
\mspage{10}
Le sous-groupe $\mathcal E_\tau$ de $\mathcal E$ image inverse de $\mathbb Z/2$ engendré par $\tau$, est le groupe cartographique ``triangulé pondéré'' de Malgoire, ayant comme gén $\tau_0,\tau_1,\tau_\infty$ \deleted{\unsure{relati.|relat.}} avec
\[
(17)\quad\left\{\begin{aligned}
\deleted{\unsure{\tau_1}\tau_\infty={}}\ l_0&=\tau_{\deleted{\unsure{0}}\infty}\tau_{\deleted{\unsure{\infty|0}}1}\\
\deleted{\tau_\infty\tau_0={}}\ l_1&=\tau_{\deleted{\unsure{\infty}}0}\tau_\infty\\
l_\infty&=\tau_1\tau_0
\end{aligned}\right.
\]
\margin{\uline{NB}\quad\begin{tikzpicture}[baseline=0.2ex]\draw (0,0.12)--(0.32,0.15);\fill (0.17,0.06) circle (0.015);\end{tikzpicture}}
$\mathcal E_\tau$ peut être considéré comme le groupe ayant comme \inserted{gén. les $\tau_i$, avec les} seules relations (16) --- on a un hom signature évident $\tau_i\longrightarrow-1$ de $\mathcal E_\tau$ dans $\pm1$, dont le noyau est $\Pi$ engendré par les $l_i$ décrits par (17)\ldots\ .

\figure{\begin{tikzpicture}\draw (0,0)--(0.8,0.025);\end{tikzpicture}}
\margin{\begin{tabular}{@{}l@{}}On a les\\relations\\[2ex]$\rho$\end{tabular}}
Revenons à une description plus symétrique des autom. extérieurs à lacets de $\mathcal E$. On désignera par $\hat{\mathcal E}$ le complété profini de $\mathcal E$, extension de $\hat\Gamma_{0,3}=\Gamma_{0,3}=\unsure{\mathfrak S_3|\mathcal S_3}\times\mathbb Z/2$ par $\hat\Pi$, et par $\widehat{\widehat{\mathcal E}}$ le groupe des autom. \deleted{\unsure{extérieurs}} \inserted{à} lacets de $\hat\Pi$, extension de $\widehat{\widehat\Gamma}=\operatorname{Aut}\operatorname{ext}_{\mathrm{lac}}(\hat\Pi)$ par $\hat\Pi$
\diagram{\begin{tikzcd}[column sep=large,row sep=small]
(18)\quad&1\arrow[r]&\hat\Pi\arrow[r]&\widehat{\widehat{\mathcal E}}\arrow[r]&\widehat{\widehat\Gamma}\arrow[r]&1\\
&&&&\unsure{\Vert|\cup}&\\
&&&&\operatorname{Aut}\operatorname{ext}_{\mathrm{lac}}(\hat\Pi)&
\end{tikzcd}}
\mspage{22}
